\chapter{Powering Up}
\label{ch:powerup}

\begin{wbwarning}
The machine begins moving on its own as soon as it is switched on. Clear the
stage and keep hands away before applying power.
\end{wbwarning}

\section{What you will see}

The configuration list appears first, then the machine reports two operations
in turn:

\begin{enumerate}[leftmargin=*]
\item \msg{HOMING} --- \msg{STARTED}. The Y stage travels to find its limit
  switch. If the switch is already made, the stage first backs off and
  approaches again, so the reference is always taken from the same direction.
\item \msg{HOMING} --- \msg{OPERATION SUCCEEDED}.
\item \msg{BONDER INIT} --- \msg{STARTED}. The bond head initialises.
\item \msg{BONDER INIT} --- \msg{OPERATION SUCCEEDED}.
\end{enumerate}

The machine is then ready. Load a configuration and you can bond.

\begin{wbnote}
Only the Y axis is homed. Z measures its height absolutely and needs no
reference, and the T axis works in relative moves throughout.
\end{wbnote}

\section{If homing fails}

Homing and initialisation are the two operations the machine treats as
critical. If either fails you will see:

\begin{lcdscreen}
\begin{lcdverb}
ERROR
HOMING
ERROR SYSTEM LOCKED
\end{lcdverb}
\end{lcdscreen}

The machine will not move and the keypad is locked. This is deliberate --- a
machine that does not know where its stage is must not travel.

\begin{enumerate}[leftmargin=*]
\item Check that nothing is obstructing the Y stage.
\item Remove anything resting on the stage or in its path.
\item Press \keycap{RESET} to restart and home again.
\end{enumerate}

If it fails repeatedly with a clear path, the limit switch or its wiring needs
service.

\section{What is remembered between sessions}

\begin{center}\small
\begin{tabular}{@{}L{50mm}L{72mm}@{}}
\toprule
\textbf{Remembered} & \textbf{Not remembered} \\
\midrule
All saved configurations & Which configuration was loaded --- the machine
                           starts at the list \\
The Settings page, once saved & Unsaved parameter edits \\
The force offset from \keycap{SETUP} & The area light on/off state \\
The tachometer offset & Clamp open/closed state \\
\bottomrule
\end{tabular}
\end{center}

\begin{wbwarning}
Parameter edits are lost at power-down unless you press \keycap{SAVE}. The same
applies to the Settings page. If you have spent time developing a recipe, save
it before switching off.
\end{wbwarning}

\section{Starting work}

\begin{enumerate}[leftmargin=*]
\item Wait for homing and initialisation to finish.
\item On the configuration list, use $\leftarrow$ and $\rightarrow$ to pick the
  bonding mode you want.
\item Move the \texttt{*} to your configuration and press \keycap{ENTER}.
\item Check the parameters on the six screens.
\item Press \keycap{TEST} to confirm the transducer is healthy.
\item Bond a test piece before starting production.
\end{enumerate}
